\subsection{Ordered vector spaces}

A preorder structure, on a vector space E, denoted by \(x \preccurlyeq y\) or \(y \succcurlyeq x\) , is compatible with the vector space structure of E if it satisfies the following two axioms;

\((\mathrm{EO}_{\mathrm{I}})\) If \(x \preccurlyeq y\) then \(x + z \preccurlyeq y + z\) for all \(z \in \mathrm{E}\) .

\((\mathrm{EO}_{\mathrm{II}})\) If \(x \succcurlyeq 0\) then \(\lambda x \succcurlyeq 0\) for every scalar \(\lambda \geqslant 0\) .

The vector space E, carrying these two structures, is called a preordered vector space (resp. an ordered vector space when the relation of preorder on E is an order).

Note that axiom \((EO_{I})\) means that the preorder structure and the additive group structure of E are compatible, that is to say, E carrying these two structures, is a preordered group (A, VI, p.~3).

Example. --- On the space \(E = R^{A}\) of all finite real-valued functions defined over A, the relation of order given by « for all \(t \in A\) , \(x(t) \leqslant y(t)\) » is compatible with the vector space structure of E.

PROPOSITION 13. --- (i) The set P, of elements \(\succcurlyeq 0\) , of a preordered vector space E, is a pointed convex cone.

\begin{enumerate}
\def\labelenumi{(\roman{enumi})}
\setcounter{enumi}{1}
\item
  Conversely, if \(\mathbf{P}\) is a pointed convex cone in \(\mathbf{E}\) , then the relation \(y - x \in \mathbf{P}\) is a preorder relation on \(\mathbf{E}\) , and the preorder structure that it defines is the only one that is compatible with the vector space structure of E and for which P is the set of elements \(\succcurlyeq 0\) .
\item
  The relation \(y - x \in \mathbf{P}\) , with \(\mathbf{P}\) a pointed convex cone, is an order relation on \(\mathbf{E}\) if and only if \(\mathbf{P}\) is a proper cone.
\item
  Axioms \((\mathrm{EO}_{\mathrm{I}})\) and \((\mathrm{EO}_{\mathrm{II}})\) imply \(P + P \subset P\) and \(\lambda P \subset P\) for all \(\lambda > 0\) . As \(0 \in P\) , it follows that P is a pointed convex cone (II, p.~11, prop. 10).
\item
  Conversely, if \(\mathbf{P}\) is a pointed convex cone, the relation \(\mathbf{P} + \mathbf{P} \subset \mathbf{P}\) implies that the relation \(y - x \in \mathbf{P}\) is a preorder compatible with the additive group structure of E (A, VI, p.~3, prop. 3); clearly writing it \(x \preccurlyeq y\) , the set \(\mathbf{P}\) is identical with the set of \(x \geqslant 0\) ; further the relation \(\lambda \mathbf{P} \subset \mathbf{P}\) for all \(\lambda \geqslant 0\) shows that axiom \((\mathrm{EO}_{\mathrm{II}})\) is satisfied. (iii) To say that \(\mathbf{P}\) is proper means that \(\mathbf{P} \cap (-\mathbf{P}) = \{0\}\) (II, p.~11, cor. 2), hence that \(y - x \in \mathbf{P}\) is an order relation.
\end{enumerate}

Example. --- * Let H be a real Hilbert space; in the vector space \(\mathcal{L}(\mathrm{H})\) of continuous endomorphisms of H, the positive hermitian endomorphisms form a proper pointed convex cone; this cone, therefore, defines an order structure compatible with the vector space structure of \(\mathcal{L}(\mathrm{H})\) and for which the relation \(A \leqslant B\) means that \(B - A\) is a positive hermitian endomorphism. *

For any pointed convex cone P in the vector space E, the set \(P \cap (-P)\) is a vector subspace, H, of E (II, p.~11, cor. 2). The canonical image \(P'\) of P in E/H is a convex cone and the inverse image of \(P'\) in E is P. Thus \(P' \cap (-P') = \{0\}\) , and \(P'\) defines an order structure on E/H that is compatible with its vector space structure.

A linear form f on a preordered vector space E is said to be positive if \(x \geqslant 0\) in E implies \(f(x) \geqslant 0\) . Or, alternatively, if the convex cone P of elements \(\geqslant 0\) in E is contained in the half space of those x for which \(f(x) \geqslant 0\) . Clearly, in the dual E* to E, the set of positive linear forms is a pointed convex cone.

